% =====================================================================
%  Professioni 2 — Il game designer e il level designer
%  Fonte: trascrizione della presentazione del corso di game design ai
%  Dev Days (AIV2): un incontro online su Zoom di circa un'ora e mezza,
%  il lunedì dalle 11:35 alle 13 circa, tenuto da due docenti di AIV:
%  Jack (Giacomo), docente di game design a Roma e online, e Andrea,
%  docente di level design. Circa 35 partecipanti di tutti i corsi del
%  primo anno, quasi tutti in webcam; domande a voce e in chat, con un
%  lungo spazio finale di domande e risposte.
%  Pacchetti richiesti nel preambolo del libro:
%    \usepackage[utf8]{inputenc}, \usepackage{booktabs},
%    \usepackage{tikz} con
%    \usetikzlibrary{arrows.meta,positioning,calc,fit}
%
%  Collocazione nel libro: parte "Alcune Professioni nel Mondo
%  Videoludico", dopo AIV1_dev_days_e_metodo_aiv (3.3).
%  Tutte le etichette di questo capitolo hanno il prefisso "p2:".
%  La figura è una ricostruzione schematica di quanto i relatori
%  spiegano a voce, non una copia delle slide.
% =====================================================================

\chapter[Il game designer e il level designer]{Il game designer e il
level designer: che cosa fanno, e che cosa vorranno da voi}
\label{p2:cap:gamedesign}

Questa è la prima delle presentazioni dei singoli corsi dei Dev Days (si
veda il capitolo~\ref{p1:cap:devdays}): quella del corso di \textbf{game
design}, che comprende anche il \textbf{level design}, ``un po' un tutt'uno''.
L'incontro vuole rispondere ``una volta per tutte'' alla domanda: \textbf{che
cosa fa un game designer?} Qual è il suo ruolo, il suo scopo all'interno di
un progetto? I relatori lo dicono subito: l'incontro sarà ``meno utile per
coloro che sono iscritti al corso di game design e più utile agli altri''.
Gli studenti di game design ``questa cosa la sperimenteranno sulla loro pelle
nell'arco di poche settimane''; agli altri serve sapere che cosa faranno i
loro colleghi, e soprattutto \textbf{che cosa un designer chiederà loro}
quando, dal secondo anno, lavoreranno insieme.

I relatori sono due docenti di AIV, che chiedono a tutti di essere chiamati
per nome e di essere trattati con il tu (``in AIV ci si dà tutti quanti del
tu, abituatevi a questa cosa; non voglio sentirmi dare del lei o chiamare
professore'').
\begin{itemize}
  \item \textbf{Jack}, ``all'anagrafe Giacomo, ma nessuno mi chiama
        Giacomo'', docente del corso di game design di Roma e online, 29
        anni; ``nella vita, oltre al docente, faccio anche il designer, e
        altre cose che piano piano scoprirete''. È lui a condurre la prima
        parte; dovrà uscire dieci o quindici minuti prima della fine per un
        imprevisto.
  \item \textbf{Andrea}, docente della parte di level design del corso; al
        di fuori dell'accademia lavora come sviluppatore indie. ``Non è lo
        stesso Andrea che avete conosciuto'' stamattina (il direttore della
        didattica): in AIV, scherza Jack, ``c'è una grande quantità di
        Andrea, ci sono due Jack e io sono uno di quelli''; due anni prima,
        in una sola classe, senza contare i docenti, c'erano cinque Andrea.
\end{itemize}

L'incontro si tiene online, su Zoom, il lunedì della prima settimana,
subito dopo l'apertura dei Dev Days; con ogni probabilità è \textbf{lunedì 5
ottobre 2026}\footnote{Per la deduzione della data si veda il
capitolo~\ref{p1:cap:devdays}. I relatori chiedono se tutti hanno
partecipato ``alla riunione di lancio di stamattina con Andrea e Daniele'',
e l'apertura dei Dev Days aveva annunciato ``alle 11.30 l'incontro per game
design''.}. Si comincia alle 11:35 ``a prescindere da quanta gente sia
entrata'' e si finisce verso l'una. I partecipanti sono una trentina abbondante,
quasi tutti con la webcam accesa.

Una sola avvertenza di lettura: la trascrizione è automatica, storpia nomi
(``Hive'' e ``Ive'' per AIV, ``City Project'' per CD Projekt Red) e in un
punto ripete per decine di volte la stessa frase (``Allora, Federico?''),
perdendo un tratto della risposta a una domanda; il capitolo lo segnala dove
accade.

Il percorso del capitolo:
\begin{enumerate}
  \item l'avvio e il censimento dei corsi (sezione~\ref{p2:sec:avvio});
  \item \textbf{che cosa fa un designer}: le risposte dei partecipanti e
        l'analogia con il regista (sezione~\ref{p2:sec:designer});
  \item il designer nelle tre fasi: pre-produzione, produzione,
        post-produzione (sezione~\ref{p2:sec:fasi});
  \item il \textbf{level design} (sezione~\ref{p2:sec:level});
  \item come cambiano i ruoli con la dimensione del team
        (sezione~\ref{p2:sec:team});
  \item la domanda centrale: \textbf{che cosa vorrà da me un designer?},
        con gli errori degli anni passati, lo scope, la fattibilità e la
        collaborazione con i programmatori (sezioni
        da~\ref{p2:sec:vorra} a~\ref{p2:sec:programmatori});
  \item il \textbf{processo iterativo} (sezione~\ref{p2:sec:iterativo});
  \item le \textbf{game jam} e il designer in una jam
        (sezione~\ref{p2:sec:jam});
  \item le domande finali (sezione~\ref{p2:sec:domande}).
\end{enumerate}

% ---------------------------------------------------------------------
\section{L'avvio}
\label{p2:sec:avvio}

Mentre i partecipanti entrano, Jack ringrazia chi ha acceso la webcam
(``fate vedere camere molto più belle della mia''), cambia il proprio sfondo
(``ora mi sento a casa'') e chiede di tenere \textbf{la webcam accesa e il
microfono spento}, accendendolo per parlare. Per prendere la parola si usa
la funzione \textbf{alza la mano} di Zoom (``non così'', ma con il tasto
\emph{Reagisci}, in basso al centro; in inglese \emph{React}).

Sulla webcam Jack è più morbido del direttore della didattica: se dà
problemi di rete, ``tenetela spenta tranquillamente''; se si riesce a tenerla
accesa senza problemi è preferibile, ma se l'alternativa è non sentire o non
riuscire a seguire, meglio seguire. In generale, però, quando si partecipa
online alle lezioni bisogna cercare di avere una webcam funzionante e
poterla accendere: se non si vede lo studente è più difficile interagire,
``vedere uno schermo nero è anche un po' brutto'', e chi non si vede attivo
\textbf{potrebbe non essere considerato presente}. (Intanto Jack nota le
case dei partecipanti: ``qualcuno in cucina, qualcuno illuminato
letteralmente dalla luce del Signore''.)

\subsection{Il censimento}

Per alzata di mano Jack fa un censimento. Molti sono di game design; alcuni
di 3D art (fra cui Giordano), alcuni di programmazione, alcuni di narrative
design (come Eleonora, che Jack prova a indovinare: ``grafica? 3D? No,
narrative''); nessuno di \emph{strategic writing}, di character animation, di
sound design o di altri corsi. ``Fa un po' ridere'', commenta: l'incontro
pensato soprattutto per i non designer è frequentato per lo più da designer.
``Questo non significa che non sarà per niente utile.'' Alcuni partecipanti
Jack li ha già conosciuti al colloquio di ammissione.

% ---------------------------------------------------------------------
\section{Che cosa fa un designer?}
\label{p2:sec:designer}

Jack comincia con una domanda, rivolta a chi non fa game design: secondo
voi, che cosa fa un designer? ``Non siate timidi, non ho potere di
bocciatura in questo momento.''
\begin{itemize}
  \item \textbf{Eleonora} (narrative): stabilisce le \textbf{regole del gioco}
        e le \textbf{meccaniche di base}, ``senza curarsi, almeno
        all'inizio, dell'aspetto estetico'', ma della parte puramente
        meccanica. Jack: interessante, ``non è sbagliata'', ma lo ha colpito
        che abbia voluto specificare che il designer non si cura
        dell'aspetto estetico; lo sottolinea, ``ci torneremo più avanti''.
  \item \textbf{Giordano} (3D art): è quello che ha la \textbf{visione ad
        alto livello} del gioco, decide i sistemi e le meccaniche ma non
        l'implementazione, che gestirà chi fa quel lavoro nello specifico;
        deve però vedere \textbf{l'integrazione di tutte le parti}.
  \item \textbf{Edoardo}: si occupa anche del \textbf{bilanciamento}, ``proprio
        a livello di numeri''. ``Assolutamente sì.''
\end{itemize}
Tutte e tre le risposte sono corrette.

\subsection{Il designer come regista}

Jack, ``amante del cinema'', propone un'analogia con la produzione di un
film: il designer è come il \textbf{regista}. Il regista non crea i costumi,
ma è il perno intorno a cui ruotano tutti i ruoli, perché ha \textbf{la
visione} di ciò che sarà il film. Così il designer: ha la visione di ciò che
sarà il gioco.

Il ruolo del designer è quindi \textbf{contemporaneamente creativo e
pratico}, e in più ha un ruolo, ``mi viene da dire'', \textbf{sociale},
perché deve connettere i vari aspetti del team. Se grafici e programmatori
possiedono soprattutto \textbf{hard skill}, abilità nei loro campi
specifici, il designer, oltre alle hard skill che deve sicuramente avere,
deve possedere anche le \textbf{soft skill} di interazione e di gestione del
team.

Attenzione, però: questo \textbf{non significa} che il designer sia più
importante, più centrale degli altri, o che abbia ``un ruolo di
comando'' (``uso questo termine un po' spaventoso apposta''). ``Il designer è
un membro del team come tutti gli altri, soltanto è il membro del team a
cui il team affida il ruolo di \textbf{guida}.''

% ---------------------------------------------------------------------
\section{Il designer nelle tre fasi}
\label{p2:sec:fasi}

A volte le persone non hanno chiaro il processo di creazione di un gioco.
Jack non entra nei dettagli (``ci saranno delle lezioni apposite''), ma
spiega che a grandi linee si sviluppa in \textbf{tre macro step}, ``che
guarda caso sono gli stessi del cinema'': pre-produzione, produzione e
post-produzione.

\subsection{Pre-produzione}

La \textbf{pre-produzione} è tutto ciò che si fa prima dell'inizio dei
lavori effettivi. Che cosa? I partecipanti rispondono:
\begin{itemize}
  \item Giordano: la \textbf{pianificazione} del progetto, gli
        \textbf{storyboard}, \textbf{trovare un finanziatore}; ``tutti e tre
        giusti'';
  \item l'\textbf{analisi di mercato}, ``giustissimo'';
  \item Sara: fare \textbf{brainstorming}, valutare pro e contro del
        progetto che sta per essere messo in pratica, cioè
        \textbf{l'analisi del proprio stesso progetto};
  \item il \textbf{concept} del gioco: ``senza quello, nessun progetto può
        partire''.
\end{itemize}
Jack aggiunge un elemento importante: oltre al proprio gioco, bisogna
\textbf{analizzare altri giochi}, per capire se il progetto ha senso. Magari
si scopre che il gioco a cui si lavora è identico a uno uscito cinque anni
prima e che è stato ``un flop clamoroso'', perché non piaceva a nessuno. A
quel punto le opzioni sono due: o si ha ``una sorta di senno di poi
antecedente'' e si \textbf{taglia} il gioco; o si capisce che cosa non ha
funzionato in quel gioco e si cerca di \textbf{non ripeterne gli errori}.
Tutto questo lo fa generalmente il designer.

\subsection{Produzione: la documentazione}

Poi c'è la \textbf{produzione}, il centro dell'esperienza di qualsiasi
membro di un team di sviluppo: ognuno lavora nel suo ambito. Ma che cosa fa,
a quel punto, il designer, oltre a sostenere e gestire il team e a occuparsi
delle meccaniche? Eleonora: produce la \textbf{documentazione} che gli altri
reparti useranno come linee guida, l'indicazione di che cosa va
implementato, come, e quanto è importante. ``Esattamente: hai detto la parola
più importante di tutte, documentazione.''

La documentazione è il centro del lavoro del designer in produzione, per due
ragioni.
\begin{enumerate}
  \item Il designer ha un'idea, ma non può comunicarla a voce a ciascuna
        persona della produzione. ``Come dicevano i romani, \emph{verba
        volant, scripta manent}'': una cosa detta a voce tra due settimane
        è dimenticata, mentre un documento scritto si può riaprire e
        rileggere.
  \item Soprattutto, è \textbf{un errore alla base} se un membro del team,
        per capire qualcosa del design, deve ricorrere direttamente al
        designer. Esempio: in uno sparatutto il designer ha deciso che
        l'arma X ha gittata Y e cadenza di fuoco Z. Se il programmatore che
        programma l'arma non trova queste informazioni e deve chiederle al
        designer, ``allora c'è un problema: il designer non ha fatto una
        buona documentazione''. Il designer deve dare a tutto il team
        accesso agli strumenti per portare avanti il proprio lavoro.
\end{enumerate}

\subsection{Post-produzione}

Infine la \textbf{post-produzione}: tutto ciò che avviene \textbf{dopo il
rilascio}. Le patch per i bug ``che immancabilmente avvengono, si presentano
e spuntano fuori solo ed esclusivamente quando si passa al contatto con
l'utenza''; guardare i feedback, capire che cosa si può migliorare,
potenzialmente creare un \textbf{DLC} (``se qualcuno non sa che cos'è un DLC,
lo dica adesso'').

% ---------------------------------------------------------------------
\section{Il level design}
\label{p2:sec:level}

Prima di andare avanti Jack passa la parola ad Andrea, per parlare del
\textbf{level design}, che ``è parte del game design, ma anche con ruoli
diversi''. La gran parte del lavoro del game designer, spiega Andrea, il
level designer la fa ``non direi tale quale'', ma con \textbf{le stesse
metodologie e gli stessi obiettivi}.

\subsection{Che cos'è un livello}

Il punto in cui game design e level design si separano è \textbf{che cosa si
produce e si idea}. Il game design si focalizza principalmente sui
\textbf{sistemi}, le \textbf{meccaniche}, gli ``ingranaggi core'' del gioco;
il level design guarda più nello specifico agli \textbf{spazi di gioco},
quelli che comunemente si chiamano \textbf{livelli}. I livelli si possono
definire in tanti modi (da giocatori, si sono visti giochi con livelli di
ogni tipo), ma hanno tutti una caratteristica comune: \textbf{sono lo spazio
nel quale avviene il gioco}, uno spazio ``chiamiamolo metafisico'' dentro il
software. Il level designer gestisce questi spazi: come sono strutturati,
cioè come il giocatore arriva dal punto A al punto B (``la definizione più
semplice''); quanto sono grandi; e soprattutto \textbf{come il giocatore
interagisce con essi}, perché è lì che avviene il gameplay di base.

\subsection{Il chi, cosa, dove, quando e perché}

Sia il game designer sia il level designer rispondono alle domande di base:
chi, cosa, dove, quando e perché. Il \textbf{game designer} si focalizza
molto sul \textbf{perché} e sul \textbf{cosa}; il \textbf{level designer}
principalmente sul \textbf{dove} e sul \textbf{come}. Per questo il level
designer ha in genere un'interazione più stretta con l'\textbf{engine}, dato
che il suo lavoro riguarda moltissimo la produzione effettiva dei livelli
durante la fase di produzione.

\subsection{Il level designer in pre-produzione: metriche e mood}

In pre-produzione il level designer fa un lavoro molto simile a quello del
game designer: ricerca, studio del proprio gioco e di altri giochi,
documentazione di preparazione. Due compiti specifici:
\begin{itemize}
  \item collabora con il game designer e i programmatori all'ideazione del
        \textbf{personaggio di gioco} dal punto di vista spaziale: quanto è
        grande, come funziona, quali sono le sue misure, cioè le
        \textbf{metriche di gioco} che riguardano il giocatore;
  \item studia gli spazi, collaborando molto da vicino con il reparto di
        \textbf{grafica}: come sarà il livello dal punto di vista
        dell'\textbf{atmosfera}, quale sarà il \textbf{mood}, qual è
        l'\textbf{obiettivo emotivo} che il livello vuole raggiungere, e
        come far sì che la parte estetica e la parte di gameplay si
        incontrino per raggiungerlo.
\end{itemize}
Come il game designer, insomma, lavora a strettissimo contatto con tutti
gli altri dipartimenti, ``ma questo è vero per entrambi i ruoli''.

\subsection{Il level designer in produzione: il blockout}

In produzione il level designer si occupa principalmente della fase di
\textbf{blockout}, o \emph{white boxing}, o \emph{gray boxing}, ``in tanti
modi'': la fase di prova delle meccaniche e delle idee, di
\textbf{iterazione di prototipi}. Qualcuno chiede (``domanda stupida'') se si
occupi attivamente anche del \emph{landscape}. ``Non esistono domande
stupide.'' Sì, è possibile che il level designer si occupi anche del
landscape, soprattutto nelle fasi iniziali di white boxing.

Che cosa vuol dire fare blockout? Fare \textbf{il livello tutto con dei cubi
grigi}, per fare le prove. Il level designer vive ``la grandissima parte''
del suo lavoro in livelli di questo tipo; chi ha visto video sul processo di
sviluppo su YouTube li ha certamente incontrati. Il suo compito è gestire
questi prototipi e assicurarsi che \textbf{funzionino}, così che il lavoro
dei grafici possa poi essere messo sopra qualcosa che funziona dal punto di
vista del gameplay.

Andrea cita un video sul processo di creazione dei livelli dell'ultimo
\emph{Spider-Man}\footnote{Con ogni probabilità \emph{Marvel's Spider-Man 2}
(Insomniac Games, 2023).}: all'inizio i palazzi sono cilindri o cubi
estremamente allungati, ``dei parallelepipedi, in realtà, grigi'', e
Spider-Man è un manichino in \emph{ragdoll} che si sposta a mezz'aria
dall'uno all'altro. Serviva a verificare che la meccanica dello
\textbf{swing} fosse abbastanza soddisfacente, ``anche perché è una delle
meccaniche più importanti del gioco''. Il level designer, insomma, porta la
documentazione propria e del game designer \textbf{dentro l'engine}, per
verificare che funzioni.

Jack aggiunge, perché l'immagine del level designer che vede ``il mondo
grigio per la maggior parte della propria esperienza'' è un po' triste, che
c'è anche la fase di \textbf{dressing}: l'abbellimento del livello, spesso
fatto in collaborazione fra level design e grafica, a seconda di come è
strutturato il team, ma che può essere anche compito del level designer.

\subsection{Il level designer in post-produzione}

In post-produzione il level designer si occupa sicuramente della fase di
\textbf{testing definitivo e feedback}, ``la parte fondamentale del lavoro di
qualunque designer''. E si occupa di concludere il lavoro sulla
\textbf{parte estetica}: la realizzazione estetica del gioco parte da
un'idea di atmosfera e di mood, e il level designer si assicura che il
lavoro dei grafici e il suo si incontrino bene nel look definitivo del
livello.

% ---------------------------------------------------------------------
\section{I ruoli cambiano con il team}
\label{p2:sec:team}

Andrea tiene a precisare una cosa che ``potrebbe non essere un'ovvietà'':
\textbf{tutti i compiti di cui si parla possono variare enormemente in base
alla grandezza del team e del progetto}. Lavorare in un indie è molto
diverso da lavorare in un tripla A.
\begin{itemize}
  \item \textbf{Più il team è piccolo} e indipendente, più il lavoro del
        designer abbraccia un po' tutto il lavoro degli altri. In un team di
        cinque persone ci sarà in genere un designer, al massimo due, che si
        occuperanno di aspetti diversi e faranno un po' di tutto.
  \item \textbf{Nei team grandi} arrivano i \textbf{producer}, i
        \textbf{lead}, i \textbf{director}, che fanno controlli ``a volo
        d'aquila'' sul resto del gioco, e il ruolo del designer si
        frammenta in tantissime sottocategorie: \textbf{combat designer},
        \textbf{economy designer}, \textbf{system designer}, \textbf{quest
        designer}, per citarne alcuni. Tutto questo è possibile solo quando
        il team conta decine, se non centinaia, di persone.
\end{itemize}
Sulla gerarchia dei team grandi si tornerà nelle domande
(sezione~\ref{p2:sec:gerarchia}).

% ---------------------------------------------------------------------
\section{La domanda centrale: che cosa vorrà da me un designer?}
\label{p2:sec:vorra}

Fatta questa premessa, si arriva alla domanda ``più importante di questa
lezione'', il suo vero scopo: \textbf{che cosa vorrà da me un designer?}

Il motivo è che, nel corso degli anni, in AIV si sviluppano videogiochi che
hanno anche la possibilità di essere pubblicati. Al primo anno il corso di
game design sviluppa un \textbf{gioco da tavolo}; dal secondo si comincia a
sviluppare giochi tutti insieme. ``Guardate bene le persone presenti in
questa chiamata, perché potreste trovarvi in team con loro tra un anno.'' A
quel punto verrà chiesto di lavorare \textbf{esattamente come un team di
produzione}, con ruoli definiti (anche ruoli interni) e con delle
\textbf{milestone}: un termine che chi non conosce forse conosce come
\emph{deadline}, e che indica \textbf{revisioni} settimanali, mensili o
bimestrali, a seconda dei producer e degli \emph{stakeholder}, in cui bisogna
dimostrare di stare effettivamente lavorando e di aver fatto progressi
rispetto alla milestone precedente.

\subsection{Tre errori degli anni passati}

Questi incontri nascono per aiutare a capire i ruoli e lavorare meglio,
perché negli anni ci sono stati problemi. Jack ne racconta tre (``e
attenzione, questa non è un'apologia dei designer e uno sfogo contro i
programmatori'').
\begin{enumerate}
  \item \textbf{Il gioco gonfiato dai programmatori.} In un team i designer
        avevano scelto un progetto; i programmatori hanno detto ``oh, figo,
        facciamolo molto più grande, perché siamo in grado di farlo''. I
        designer hanno detto di no; i programmatori l'hanno fatto lo stesso,
        fra di loro. ``Volete provare a indovinare com'è andata? Molto male.
        Non c'era un gioco, alla fine.''
  \item \textbf{Il mondo grigio.} Un gruppo di design credeva che i propri
        grafici fossero in grado di produrre asset ``iper complessi ed
        estremamente dettagliati'' in giorni e non, com'era in realtà, in
        mesi. Quando il gioco è stato consegnato era ``letteralmente il mondo
        grigio e tetro che aveva descritto Andrea, con due cose messe qua e
        là'', perché i designer non erano riusciti a intuire le capacità dei
        grafici con cui lavoravano: ``errore loro''.
  \item \textbf{La narrativa scollegata.} Game designer e narrative designer
        non si erano parlati, e il gioco andava in una direzione e la
        narrativa in un'altra. Il gioco permetteva, a livello meccanico, di
        \textbf{viaggiare fra linee temporali} per risolvere puzzle
        ambientali; la narrativa parlava di \textbf{una guerra con gli
        alieni}. ``Per carità, magari ha senso'', ma non c'era nessun nesso
        fra le due cose: ``ehi, aspetta, c'è quest'albero caduto, se vado nel
        passato dovrei poter trovare un punto in cui l'albero è ancora in
        piedi e superare l'ostacolo''; e intanto il dialogo diceva ``ci
        stanno attaccando, arrivano dal cielo, come faremo a salvarci?''.
\end{enumerate}

\subsection{Che cosa bisogna sapere}

Dal punto di vista del game design, quindi, bisogna anzitutto sapere ciò che
si è detto: che cosa fa un designer in pre-produzione e in produzione (la
post-produzione ``si vedrà più avanti''). Ma ogni parte del team ha il suo
ambito.

\paragraph{Game designer e narrative designer.} È molto importante che si
parlino \textbf{costantemente}, sia per sviluppare una storia coerente con
l'idea del gioco, sia per dare un senso alle singole meccaniche. I narrative
designer hanno anche un compito importante nel \textbf{comunicare il loro
concept narrativo} ai game designer, perché così si gestisce meglio anche
l'aspetto grafico. Jack racconta di un team in cui la narrativa prevedeva un
mondo all'età della pietra, \textbf{senza l'uso del metallo} per gli
strumenti. È stato fondamentale che lui, come game designer, lo dicesse ai
grafici, che stavano già cominciando a ideare utensili di metallo e hanno
dovuto ritexturizzarli per farli sembrare di osso o di legno: per la
narrativa scritta dai narrative designer era importante che il mondo fosse
primordiale.

\subsection{Comunicare e conoscere le proprie skill}

Andrea generalizza: la cosa fondamentale, per chi si deve interfacciare con
i designer, è \textbf{comunicare sempre} e tenere presenti i propri
\emph{skill set} (punti di forza, conoscenze), riconoscere se l'altra persona
li ha e, se non li ha, aiutarla a capire il proprio punto di vista, ai fini
dell'obiettivo comune: lo sviluppo del gioco. Nella situazione accademica ci
si troverà spesso ad avere i \textbf{designer come mediatori}.

Ci si troverà anche ad avere \textbf{designer come producer}. Come team di
sviluppo si risponderà a dei \textbf{lead} e a degli \textbf{stakeholder},
che sono AIV stessa, i docenti, a cui bisogna portare il gioco; e spesso
saranno i designer a organizzare le fasi di lavoro, perché il loro è il corso
che si occupa di più della parte di \emph{producing}.

Detto questo, un designer \textbf{non ha lo stesso skill set} di un
programmatore né di un grafico: può avere idee e avere bisogno del parere
degli altri sulla loro \textbf{fattibilità} e sul \textbf{metodo} con cui
svilupparle. Se un designer vuole un gioco con tre personaggi diversi e i
grafici dicono ``è impossibile farti tre personaggi umanoidi in un mese e
mezzo'', bisogna rivedere lo \textbf{scope} del progetto.

\subsection{Lo scope}
\label{p2:sec:scope}

\textbf{Scope} è ``una parola importantissima'', con cui i designer hanno a
che fare quotidianamente: sono \textbf{le dimensioni del progetto},
letteralmente le dimensioni a cui si punta ad arrivare, quanto è grande il
gioco. I designer vengono istruiti, ``con più o meno forza durante gli anni''
(``ci aspettiamo la gravità prima ancora dell'inizio del corso'', scherza
Jack), a \textbf{riconoscere la grandezza di un progetto} e quindi la sua
\textbf{fattibilità}. Ed è compito degli altri, ciascuno per il proprio
ambito di competenza, aiutare i designer a capire quanto il proprio lavoro
rientri in quelle dimensioni.

D'altro canto, i designer \textbf{non sono i capi} del progetto e non devono
diventare ``i vostri boss, anzi, per niente'': svolgono una funzione di
\textbf{mediatori}, al fine di arrivare a un progetto finito. Questo non
vuol dire che gli altri non possano partecipare all'ideazione: possono,
eccome. Ma bisogna riconoscere che \textbf{la skill principale del
designer}, o una delle principali, è \textbf{convertire l'idea in un progetto
realizzabile}.

\subsection{Il volo libero: un'idea che sembra semplice}

Andrea fa un esempio: ``vogliamo fare un gioco dove il personaggio
principale vola'', con volo 3D libero, come un uccello, su, giù, destra,
sinistra. Comporta difficoltà importanti, e potrebbe essere troppo grande per
ciò che si deve raggiungere, diciamo entro la fine dell'anno, entro i sei
mesi di lavoro (i progetti in collaborazione saranno ``più o meno di questa
grandezza'', ``a meno che non passino all'anno successivo''). Le idee di chi
non è designer sono validissime; ma c'è una figura nel team il cui lavoro è
proprio \textbf{filtrare le idee in feature attuabili} nel gioco: il
designer, game o level che sia, fa da \textbf{filtro} e da
\textbf{adattatore}, dalle idee al gioco. È una collaborazione ``a doppia
fase'' fra designer e programmatori, grafici, narrative designer, in cui
ciascuno usa le proprie hard skill per sopperire a quelle che mancano
all'altro.

\subsection{Tutti hanno un'idea per un gioco}

``Alzi la mano chi \textbf{non} crede di avere una buona idea per un gioco.''
Una persona su trenta (``grazie, Sara, per esserti esposta''). È
perfettamente normale: chiunque, per strada, alla domanda ``hai un'idea per
un gioco?'' risponderebbe di sì. Sono tutte idee valide? ``Assolutamente
no'': la maggior parte non supererebbe neanche la pre-produzione. Il problema
è che le persone non lo sanno, perché non hanno le competenze per saperlo;
bisogna fare un percorso di studi, quello che faranno i game designer.

Jack ci mette ``non una, ma entrambe le mani sul fuoco'': entro la fine
dell'anno tutti gli studenti di game design andranno da lui a dirgli ``Jack,
avevo un'idea per un gioco, adesso ho capito quanto poco fosse fattibile,
adesso ce n'è una migliore''. Lui risponderà: ``aspettate il secondo anno, e
succederà di nuovo la stessa cosa''. È normale.

Il volo, appunto, sembra fattibile: ``se si può programmare, perché non si
può fare?''. Perché \textbf{stravolge tutto il level design}: non si può più
fare un level design lineare, non si può incanalare il giocatore da qualche
parte, a meno di costruire barriere ``che arrivino dal cielo alla terra''. A
quel punto la domanda da farsi è: \textbf{perché} stiamo facendo un gioco in
cui il personaggio vola? Si pensi a quanto possono cambiare le dimensioni e
il level design di un progetto anche solo passando dal 2D al 3D: se succede a
lavori in corso, ``buttate tutto'', il 100\% del lavoro, ``all caps''. Oppure
si tenta ``una follia'', un gioco a volte 2D e a volte 3D, che si rivelerà
``probabilmente la zappata sui piedi più grande che abbiate mai
sperimentato''.

Per essere ``molto, molto chiari'': questo non vuol dire che le idee non
siano valide, ma che \textbf{vanno validate}. Magari c'è qualcosa di
funzionale in un'idea; ma il momento in cui l'idea si confronta con la realtà
del progetto è quando passa nelle mani del designer, che usa le sue hard
skill per ``macinarla'' e vedere quanto è fattibile.

\subsection{Il reality check, nei due sensi}

Vale anche il contrario: \textbf{bisogna aspettarsi che i designer vengano a
chiedere dei reality check}, verifiche sulla fattibilità. Se un designer
chiede ``un livello tipo \emph{Scorn}'' (``tutto horror ambientale'', il cui
stile visivo è lo stesso di \emph{Alien}\footnote{Più precisamente, lo stile
visivo di \emph{Scorn} (Ebb Software, 2022) si ispira esplicitamente
all'arte di H.~R.~Giger, l'artista che disegnò la creatura di \emph{Alien}
(1979); non è ``lo stesso''.}, ``e si vede''), e i grafici sono in grado di
fare due asset di quel tipo ma decisamente non un intero livello modulare, è
``legittimissimo'' rispondergli ``cosa stai dicendo? No, decisamente no''.
(``Potresti anche picchiarlo.'' ``Non è vero, non incentiviamo la violenza'',
ma sarebbe una richiesta così al di là della fattibilità che una reazione
violenta sarebbe ``non dico accettabile, ma quantomeno comprensibile''.)

Altri esempi di Jack: se chiedesse una narrativa con almeno
\textbf{diecimila finali alternativi}, ciascuno con le sue cutscene, i
narrative designer avrebbero tutte le ragioni di prendersela; ai
programmatori non può chiedere di programmare \textbf{venti IA nemiche},
ognuna con un comportamento indipendente ma in comunicazione con le altre
per creare sistemi complessi. ``O meglio, posso, se ho dieci milioni di
budget e un centinaio di persone nel team.'' Sono, precisa Andrea,
\emph{reductio ad absurdum}: esempi portati all'eccesso, soprattutto rispetto
allo \textbf{sviluppo indie}, che sarà il modello dei progetti in accademia.
I team di AIV lavoreranno molto più come team indie che come team doppia o
tripla A, perché non ne hanno i numeri ``e, spoiler, non avete i fondi che
hanno i team doppia o tripla A''.

% ---------------------------------------------------------------------
\section{Designer e programmatori: non a compartimenti stagni}
\label{p2:sec:programmatori}

Andrea torna su una delle risposte iniziali, quella di Giordano: i designer
si occupano della struttura delle meccaniche, ma non del metodo di
implementazione. ``Questa cosa è vera fino a un certo punto.'' A essere
fiscali, il designer struttura \textbf{la funzione}, come la meccanica
dovrebbe funzionare nel gameplay, e il programmatore pensa a \textbf{come}
farla funzionare così; in teoria le due cose si potrebbero separare. Nella
realtà, spesso designer e programmatore discutono direttamente anche del
metodo di realizzazione. Andrea lo porta dalla propria esperienza: il
programmatore legge il documento e chiede ``ok, ma come vuoi che funzioni
esattamente? Perché ho almeno cinque metodi diversi per farla. Come vuoi che
la faccia?''. È un buon esempio del \textbf{processo iterativo} della
produzione. Il designer deve quindi immaginare non solo la funzione della
meccanica nel gioco, ma anche dare \textbf{dei punti fermi
dell'implementazione} (``vorrei che questa cosa avesse almeno x, y e z, fatte
in questo modo''): una linea guida per il programmatore, che dirà ``questo ha
senso, questo no'', in un lavoro di collaborazione.

\subsection{Un gioco fatto a compartimenti stagni}

Una delle cose fondamentali dello sviluppo videoludico è che \textbf{non è
un lavoro a compartimenti stagni}, ``assolutamente no'': è un lavoro
``estremamente fluido''. O meglio: c'è chi ha provato a farlo a compartimenti
stagni, ``e non è andata molto bene''. Andrea propone un indovinello: un
gioco di una casa molto famosa, rilasciato ufficialmente ``a fine 2019''
dopo vari ritardi, e talmente buggato al rilascio da far quasi fallire la
casa che l'ha prodotto. ``Cyberpunk, esatto.''

\emph{Cyberpunk 2077}, dice Andrea, è stato prodotto, ``per ammissione stessa
del team di CD Projekt'', con una metodologia chiamata \textbf{waterfall},
cioè \textbf{a compartimenti stagni}, che è in sostanza come funziona la
produzione cinematografica\footnote{Due precisazioni. \emph{Cyberpunk 2077}
è uscito il 10 dicembre 2020, non a fine 2019. Quanto alla metodologia
``waterfall'', è la ricostruzione del relatore; il capitolo non ha modo di
verificare la dichiarazione dello studio a cui si riferisce.}. E il
risultato è stato un gioco che, per arrivare alla sua vera forma definitiva,
ci ha messo ``più o meno altri due o tre anni'': la versione 2.0 ha
realizzato il vero obiettivo del gioco. ``È un giocone, ragazzi'', aggiunge
Jack: chi lo compra e lo gioca adesso lo faccia, ``è un gioco incredibile'';
all'epoca, purtroppo, non era in grado di presentarsi come tale.

% ---------------------------------------------------------------------
\section{Il processo iterativo}
\label{p2:sec:iterativo}

``Vi ho un po' mentito all'inizio'', ammette Jack. Il processo è
pre-produzione, produzione, post-produzione, ``ma con un grosso ma'':
finita la pre-produzione \textbf{non si passa necessariamente alla
produzione}; si potrebbe dover ricominciare la pre-produzione se si trovano
dati che non vanno bene. E in produzione questo diventa ``mille volte più
vero'': non bisogna aspettarsi che la produzione sia un unico ciclo, ci
saranno molti tentativi in cui si fa, si testa, si annotano gli errori e si
ricomincia da capo. ``Questo vale per qualsiasi cosa.''

(``Ovviamente, ragazzi, potevate alzare la mano in qualsiasi momento per fare
domande. So che è un po' brutto dirlo a posteriori, però speravo che fosse
implicito.'' Si lasceranno quindici o venti minuti alla fine solo per le
domande, perché ``non ci sono domande stupide'': i concetti, soprattutto in
game design, spesso non sono ovvi. Jack si è reso conto, lavorando come
designer, che anche persone con una certa esperienza spesso non sanno in che
cosa consista davvero il lavoro del designer.)

Andrea aggiunge: il processo iterativo ``vi colpirà tutti'', e chi non ha
mai fatto un lavoro del genere dovrà imparare a conviverci. È una struttura
di lavoro che prevede \textbf{fasi ripetitive}. Riprendendo l'esempio del
volo: il designer scrive la documentazione; la feature viene implementata;
il level designer la mette nel livello; la si valuta (che cosa va, che cosa
non va); si torna alla \emph{whiteboard}, all'implementazione o alla
documentazione; si capisce che cosa non è andato e si riprova, finché non
funziona. Lo sviluppo di un gioco, quindi, \textbf{non è una linea dritta}
dall'inizio alla fine, ma ``una specie di \textbf{fusillo}'' che continua a
girare su se stesso in vari punti; e quei punti in cui gira su se stesso
sono le \textbf{feature} del gioco, le meccaniche, le strutture di gameplay,
tante o poche a seconda della grandezza del gioco (``nei giochi indie,
tendenzialmente, saranno meno''). In game design si lavora in modo iterativo
costantemente, e lo si insegna agli studenti: ``è tutto quanto incentrato su
questo'' (figura~\ref{p2:fig:fusillo}).

\begin{figure}[htbp]
\centering
\begin{tikzpicture}[>=Stealth, font=\small]
  \draw[->, thick] (0,0) -- (1.2,0)
    .. controls (2.2,0) and (2.2,1.1) .. (1.7,1.1)
    .. controls (1.2,1.1) and (1.2,0) .. (2.4,0) -- (3.4,0)
    .. controls (4.4,0) and (4.4,1.1) .. (3.9,1.1)
    .. controls (3.4,1.1) and (3.4,0) .. (4.6,0)
    .. controls (5.6,0) and (5.6,1.1) .. (5.1,1.1)
    .. controls (4.6,1.1) and (4.6,0) .. (5.8,0) -- (6.8,0)
    .. controls (7.8,0) and (7.8,1.1) .. (7.3,1.1)
    .. controls (6.8,1.1) and (6.8,0) .. (8,0) -- (9.4,0);
  \node[below, font=\scriptsize] at (0.2,-0.05) {inizio};
  \node[below, font=\scriptsize] at (9.2,-0.05) {gioco finito};
  \node[above, font=\scriptsize] at (1.7,1.15) {feature 1};
  \node[above, font=\scriptsize, align=center] at (4.5,1.15) {feature 2\\(più giri)};
  \node[above, font=\scriptsize] at (7.3,1.15) {feature 3};
  \node[draw, rounded corners, font=\scriptsize, align=left, anchor=north]
    at (4.7,-0.7) {ogni giro: documentazione $\to$ implementazione $\to$
    inserimento nel livello\\ $\to$ valutazione $\to$ ritorno alla
    whiteboard, finché non funziona};
\end{tikzpicture}
\caption{Lo sviluppo come ``fusillo'': non una linea retta, ma un percorso
che gira su se stesso in corrispondenza di ciascuna feature. Ricostruzione
schematica dell'immagine usata a voce dal relatore.}
\label{p2:fig:fusillo}
\end{figure}

% ---------------------------------------------------------------------
\section{Le game jam e il designer in una jam}
\label{p2:sec:jam}

Federico chiede in chat: \textbf{nelle jam, quanto è importante un game
designer che sa più o meno programmare giochi?} Prima di rispondere, Jack
chiede chi non sa che cos'è una game jam; qualcuno alza la mano (``hanno
aspettato che Gabriel alzasse la mano per farsi avanti'').

\subsection{Che cos'è una game jam}

Una \textbf{jam} è un evento che non riguarda solo i videogiochi: si
conoscono le \emph{jam session} musicali, e ce ne sono nel cinema, nella
moda, nella scrittura. L'idea è una \textbf{competizione amichevole} in cui
si simula il processo di creazione in un tempo brevissimo: molto spesso
24--48 ore, a volte addirittura filate. Jack ne ha fatta una sola di 48 ore
filate: ``è un'ammazzata, è bellissimo, ma non dormire per 48 ore, o dormire
appoggiato sul computer perché sei crollato, è un'esperienza molto
particolare''.

Come funziona: all'inizio si formano i team e viene dato un \textbf{tema},
che si scopre solo a jam iniziata, così che nessuno possa arrivare con
qualcosa di già preparato. Poi bisogna realizzare in modo estremamente
rapido un videogioco ispirato a quel tema, in 12, 24, 48 ore, o nel tempo
previsto. La jam più grande del mondo è la \textbf{Global Game Jam}, che
negli ultimi anni dura, ``mi sembra'', cinque giorni, una settimana
lavorativa dal lunedì al venerdì, ``e lì si suppone che voi dormiate''
\footnote{Tradizionalmente la Global Game Jam si svolge in 48 ore, in un fine
settimana di fine gennaio; nelle edizioni più recenti l'organizzazione ha
concesso alle sedi calendari più flessibili, distribuiti su più giorni. La
durata di cinque giorni è quella ricordata dal relatore, che la dice con
qualche incertezza.}.

L'obiettivo pratico è \textbf{arrivare alla fine con qualcosa di
giocabile}: un pezzo di gioco o un gioco potenzialmente intero. L'obiettivo
umano è mettersi alla prova, conoscere persone nuove, lavorare con altri
dipartimenti, allenarsi nel game development: ``un'ottima chance per fare
pratica''. Jack consiglia ``tantissimo'' le jam: ce n'è stata una in AIV
proprio nel fine settimana appena passato, e chi non ha partecipato potrà
farlo negli anni successivi, o cercarne online, ``ce ne sono un sacco''. È
uno dei modi migliori per crescere in fretta come designer e per aggiungere
progetti al proprio portfolio.

\subsection{A che cosa serve un designer in una jam}

Ora la domanda di Federico. Jack capisce da dove nasce: si è detto che il
designer scrive documenti e analizza giochi, ma in una jam non c'è tempo.
I documenti non si scrivono ``tanto per'', si scrivono se sono utili: ``il
tempo che io ho finito di scrivere un documento per una jam, la jam è
finita''. E siccome in una jam si lavora tutti insieme, anche in cuffia, la
scrittura di un documento diventa meno utile, non per forza del tutto
inutile: in certi casi ha senso farlo in forma abbreviata, non un documento
di produzione ma un documento rapido per dare l'idea a tutti. A che cosa
serve, allora, il designer?
\begin{itemize}
  \item anzitutto per la \textbf{fattibilità}: la prima fase di una jam è
        scoprire il tema, e il designer è il controllo sul ``riusciamo
        davvero a fare in 24 ore un gioco così?'';
  \item poi, nel caso ``tra mille virgolette peggiore'' di un designer che
        sa fare solo design, senza competenze di programmazione, narrativa o
        grafica (``rarissimo che sia effettivamente così''), vale ``il bello
        della jam'': \textbf{non si può restare con le mani in mano}. Se non
        si sa fare qualcosa, \textbf{la si impara}, lì per lì. Nella sua
        prima jam Jack ha imparato a fare la \textbf{UI}, l'interfaccia
        utente, che non aveva mai studiato: aveva finito la sua parte (in
        quella jam faceva anche da narrative, aveva scritto tutti i
        dialoghi), mancavano cinque ore; un tutorial di YouTube su uno
        schermo e il lavoro sull'altro. ``Perché quando non hai idea di che
        cosa fare, chiedi che cosa manca e ti ingegni per farlo'';
  \item e se non manca niente, il designer deve comunque trovare un modo di
        \textbf{supportare il team}: una cosa che Jack ha fatto, ``che può
        sembrare la più banale del mondo ma è stata molto utile'', è stata
        andare a prendere il caffè per tutti.
\end{itemize}

Andrea, che era con Jack in quella jam (``abbiamo fatto insieme quella jam,
l'abbiamo anche vinta''), aggiunge due cose. Primo: è importante avere una
persona in più in una jam, qualunque sia il suo skill set, perché più si è,
meglio si divide il lavoro. Secondo, con una contro-domanda: ``una jam è
l'occasione per prendere un'idea di partenza e convertirla in un gioco.
Quanto è importante il ruolo di una persona il cui intero lavoro è
convertire le idee in funzioni di gioco?''. Trasformare una parola, ``non lo
so, \emph{goccia}'' (``un tema di jam a caso''), in un gioco. E all'inizio
della jam si fa \textbf{brainstorming}: i designer sono addestrati alle
tecniche di brainstorming, che non è ``spariamo le nostre idee'' ma ha
tecniche specifiche di gestione; il designer le conosce e può guidare il
brainstorming fino a un'idea.

È vero, ammette Andrea, che il game designer lavora di più nella prima metà
della jam che nella seconda; chi finisce il suo lavoro fa ciò che serve.
Soprattutto il level designer, che lavorando spesso in engine si trova a fare
blockout, livelli di prova, test delle meccaniche (``chi fa level design sa
fare game design, e spesso viceversa''). Anche il game designer può testare
ciò che gli passano i programmatori: ``ho fatto questa feature; va bene,
passami il file di gioco, la provo e vediamo se funziona, intanto vai
avanti''. È raro che un designer non sappia davvero implementare nulla, a
meno che non abbia appena cominciato, ``e ci sta''.

Andrea racconta di un amico medico che voleva vedere ``com'è il vostro mondo
lavorativo'' e ha partecipato a una jam: ``ha fatto quello che andava
fatto''. Ha montato il trailer, ha aperto la pagina itch, ha dato le sue
idee. Se una persona senza alcuna competenza, che non aveva mai fatto un
corso, può essere un membro utile di una jam, ``figuriamoci un designer, che
le competenze ne ha da spalare''.

% ---------------------------------------------------------------------
\section{Le domande finali}
\label{p2:sec:domande}

La lezione vera e propria è finita, ma restano un quarto d'ora o venti
minuti di domande: ``se c'è qualcosa che non avete capito, siate
tranquillissimi nel dirlo apertamente''.

\subsection{Come si vince una jam?}

Emanuele chiede come si vince una jam: c'è una giuria di esperti? Prima del
come, Jack chiarisce \textbf{che cosa} si vince: nella stragrande maggioranza
dei casi, \textbf{niente, se non la gloria}. Non ci sono premi in palio,
perché rischierebbero di rendere la competizione meno amichevole; il premio
è sentirsi dire ``hai fatto il gioco migliore'', e fine. In rari casi ci sono
premi, in genere non in denaro: una sponsorizzazione del gioco per
produrlo su scala più ampia, o semplicemente il bollino della jam sul gioco
e la comunicazione della vittoria. ``Si vince la gloria, e forse un po' di
fama.''

Il come varia: in alcune jam vota una \textbf{giuria}; in altre votano
\textbf{i jammer stessi}, che si testano i giochi a vicenda e indicano su un
foglio digitale il gioco preferito (non il proprio), e vince quello con più
voti; a volte è un mix. La \textbf{Firestarter} di AIV\footnote{Nel
capitolo~\ref{p1:cap:devdays} la stessa jam di inizio anno è chiamata
``Fast Starter''; il nome esatto non è ricostruibile con certezza.}, ricorda
Andrea, è un mix: c'è una giuria di docenti che testano i giochi, con un
docente per dipartimento che ha l'ultima parola, e c'è la votazione dei
jammer, tramite un modulo Google.

\subsection{Il gioco da tavolo del primo anno e i gruppi}

Un'altra studentessa fa due domande. Primo: al primo anno si realizzerà un
gioco da tavolo; quali saranno le task e gli obiettivi intermedi? Si
lavorerà in piccoli gruppi dentro il singolo corso? Secondo: dal secondo
anno, i gruppi per il videogioco li formeranno i docenti?

Alla seconda: \textbf{sì}, i gruppi li creano i docenti, anche al primo
anno, e al secondo li creano i docenti dei vari corsi. Alla prima: si
lavorerà in \textbf{piccoli gruppi all'interno del corso}, con
\textbf{milestone intermedie} che consisteranno nel \textbf{testing} del gioco
al livello in cui è, dopo il quale i docenti daranno task del tipo
``migliora questa cosa'', ``questa non va bene, cambiala'', ``questa va
bene, tienila''; e il gioco verrà perfezionato piano piano con il processo
iterativo. Le task e gli obiettivi intermedi sono \textbf{gli stessi di un
videogioco}: non a livello granulare (non si possono chiedere in un gioco da
tavolo le stesse cose di un videogioco e viceversa), ma come tipologia di
richieste e di obiettivi.

Perché al primo anno un gioco da tavolo e non un videogioco? Non perché a
AIV piacciano particolarmente i giochi da tavolo (``anche se a noi due,
specificamente, piacciono un botto''). Il motivo è che \textbf{un gioco da
tavolo richiede le stesse competenze di design di un videogioco, ma senza le
competenze tecniche degli altri dipartimenti}: non serve programmazione, non
serve grafica. Un designer, senza altre competenze, può farlo da solo:
servono ``potenzialmente pezzi di cartoncino e delle penne''.

\subsection{Come si capisce che un'idea è sbagliata?}

Tancredi chiede come si capisce che l'idea che si sta portando avanti non è
quella giusta e che bisogna ricominciare da capo: ci sono campanelli
d'allarme? ``C'è una lezione di tre ore su questo'', risponde Andrea, che ne
dà una versione breve. Tre cose.
\begin{enumerate}
  \item Un gioco parte da una \textbf{visione}: una vision globale, una
        missione che il gioco deve raggiungere da vari punti di vista.
        Quando si valuta un'idea bisogna chiedersi se rientra nella visione
        del gioco. Esempio volutamente esagerato: se si sta facendo
        \emph{Animal Crossing}, ha senso metterci dentro il \emph{parry} alla
        \emph{Sekiro}\footnote{La trascrizione dice ``il Perry alla
        Seichiro'': con ogni probabilità il sistema di deflessione degli
        attacchi di \emph{Sekiro: Shadows Die Twice}.}? ``Probabilmente no.''
  \item L'\textbf{esperienza}: guadagnandone abbastanza, i campanelli
        d'allarme ``escono fuori''. Si comincia a capire quali tipi di
        meccaniche (le idee spesso si traducono in meccaniche o in funzioni
        di gameplay) creano il \textbf{ripple effect} più grande, l'effetto
        dell'onda quando si tira un sassolino nell'acqua: alcune sono
        piccole, altre a occhio sembrano facili e invece hanno conseguenze
        molto più grandi su tutto il resto del progetto.
  \item Il \textbf{processo iterativo} stesso risponde alla domanda: se si
        implementa un'idea e il risultato non è quello voluto, ``quello è
        già il campanello d'allarme più grande'': se non funziona, bisogna
        cambiare.
\end{enumerate}

\subsection{E se nel team manca un ruolo?}

Gabriel chiede se sia capitato che nei team mancassero persone per certi
ruoli, come animatori o effetti sonori. Sì, ``ora molto meno'': negli ultimi
anni, soprattutto dal secondo anno, si cerca di formare team il più
completi possibile; ma alcuni corsi, come animazione e sound design, hanno
pochissimi studenti da distribuire. In compenso c'è ``una figata'' che AIV
permette dal secondo anno: \textbf{chiamare persone esterne} ad aiutare.
Lavoreranno su un gioco di AIV e ``non riceveranno una lira'', ma se a loro
sta bene possono partecipare: ``vieni con me, mi fai il rigging di questa
cosa?''. Va comunque concordato con la direzione.

\subsection{Chi inventa ambientazione e personaggi?}

Un'altra domanda: l'ambientazione e il personaggio li formula il narrative
designer? ``In realtà non è detto.'' Jack, game designer che si è dedicato
in alcuni momenti anche al narrative design, dice che succede spesso una via
di mezzo: quando un designer idea un gioco è molto difficile che lo immagini
soltanto con le sue meccaniche, senza un contesto narrativo. Gli è capitato
spesso di dire ``mi sono immaginato questo gioco, in un'ambientazione più o
meno così, con un protagonista più o meno così''; i narrative devono poi
dirgli se ha senso e prendere quel concept per ampliarlo in un mondo, una
storia, uno o più personaggi. È molto difficile che qualcuno abbia un'idea
che riguardi soltanto il proprio ambito: per questo si fa brainstorming
tutti insieme e si gestisce l'idea in modo da ``avere i pezzi migliori delle
idee di tutti''.

\subsection{Collaboratori da altri anni di AIV?}

Filippo chiede se i collaboratori esterni possano essere studenti di AIV di
altri anni. Possono, ``nel senso stretto del termine'', ma ogni studente ha
il progetto del suo anno: è difficile che uno del terzo anno, con il suo
progetto, venga a dare una mano a un altro. ``Per ora non è successo.''
Qualche forma di comunicazione fra anni c'è stata: per esempio studenti del
primo o del secondo anno hanno fatto da \textbf{tester} per i giochi degli
anni successivi quando serviva. Ma le collaborazioni vere e proprie fra anni
sono molto difficili; i collaboratori, di solito, sono davvero esterni.

\subsection{La gerarchia nei team grandi}
\label{p2:sec:gerarchia}

Federico fa una domanda sulla gerarchia dei designer nei team grandi. Qui la
trascrizione ripete per decine di volte ``Allora, Federico?'' e perde la
domanda e l'inizio della risposta\footnote{Dalla parte di risposta che
sopravvive si capisce che la domanda riguardava come sono organizzati,
gerarchicamente, i designer nelle grandi aziende.}. La risposta riprende
così: va immaginata come \textbf{una piramide}. In cima ci sono i
\textbf{lead}, che hanno una visione di tutto, ma molto generica; scendendo
si vede meno, ma più nello specifico: \textbf{senior} e \textbf{junior}, e
ogni ambito ha il suo specialista. Magari ci sono tre persone, o anche una
sola, che si occupano esclusivamente di \textbf{combat design}; faranno
riferimento a qualcuno che si occupa del \textbf{system design}, che farà
riferimento a qualcuno che si occupa del \textbf{game design}.

``La mia risposta preferita, Federico, è: \emph{dipende}.'' Non è uno
scherzo: non c'è una definizione universale, e cambia moltissimo a seconda
dell'azienda, del progetto, delle sue richieste. In un gioco molto incentrato
sul combattimento ci sarà probabilmente un \textbf{lead combat designer} che
gestisce un team di combat designer e combat system designer, con dentro
probabilmente anche gli \textbf{enemy designer}. Scendendo verso le singole
persone, ciascuna avrà una serie di task settimanali. Esempio: ``questa
settimana devo fare l'enemy design dei non morti base del Borgo dei Non Morti
di \emph{Dark Souls}''. Uno fa il cavaliere, uno quello con la lancia, uno
quello con l'arco; a fine settimana, chiamata con il lead combat designer per
mostrarglieli; poi passano ai programmatori. In generale, più l'azienda è
grande più tende ad assumere: aziende come \textbf{Ubisoft} fanno del numero
di dipendenti la loro forza, ``hanno migliaia di designer assunti''.

\subsection{Dove si vedono i giochi da tavolo degli anni passati?}

Giordano chiede se i giochi da tavolo del primo anno vengano esposti da
qualche parte. Va chiesto a \textbf{Daniele Fusetto}, che porta avanti i
progetti dopo il primo anno, o alla direzione di AIV, che sa dove si possono
vedere i vecchi progetti. Non tutti sono visibili: solo quelli \textbf{portati
avanti}. Alla fine del primo anno, ma anche degli anni successivi, c'è una
sorta di \textbf{revisione finale} in cui i docenti danno un
\textbf{``passato'' o ``non passato''}, che non vuol dire promosso o
bocciato: semplicemente il progetto può essere presentato a Fusetto oppure
non è abbastanza buono. Ci sono stati anni in cui sono passati tutti, anni in
cui ne sono passati uno o due, anni in cui non ne è passato nessuno. E i
progetti passati vengono presentati solo se gli studenti \textbf{vogliono}
dedicarcisi; alcuni dicono ``siamo arrivati fin qui, ma non ci interessa
andare avanti''. Alcuni giochi da tavolo di AIV sono stati portati allo stand
di AIV al \textbf{Play}, il festival del gioco\footnote{La trascrizione dice
``al Play di Bologna''. Il festival del gioco \emph{Play} si tiene a Modena;
è probabile che si tratti di quello.}, per essere testati dal pubblico,
``mi pare negli ultimi due anni''.

Filippo chiede se quindi anche il gioco da tavolo possa continuare dopo il
primo anno. Potenzialmente sì, se ha il ``passato''. Finché si è dentro
l'anno di corso, il passaggio o meno del progetto è nelle mani dei docenti;
poi passa alla direzione dei giochi da tavolo di AIV, cioè a Fusetto, che
decide se e come continua; e può capitare che un progetto che per i docenti
passa non passi a livello di direzione.

\subsection{Chi coordina: game designer o narrative designer?}

``È l'una'': ultima domanda, di Edoardo. Quanto del coordinamento fra le
varie figure è fatto dal game designer e quanto dal narrative designer?
Jack dà due risposte. Nel \textbf{mondo del lavoro}, il coordinamento non lo
fa nessuno dei due: lo fa il \textbf{producer}, una figura apposita che
gestisce il lavoro degli altri. \textbf{In AIV}, però, sia il corso di game
design sia quello di narrative hanno basi di producing, e il ruolo verrà
assegnato parzialmente a entrambi. Storicamente, negli anni passati, il
producing lo hanno fatto spesso i game designer, ma ci sono stati anche
narrative che hanno preso quel ruolo (nell'ultimo anno, per esempio, nel
team di un gioco\footnote{La trascrizione dice ``nel team Dozen'': il nome
non è ricostruibile con certezza.}, era molto dei narrative). È lasciato a
chi ha le skill e vuole prendersene cura: qualcuno dirà ``farò io il producer
per questo pezzo di progetto'', ``un game o un narrative''. Se ne parlerà.

Bisogna anche scegliere uno o più \textbf{lead}, e non è detto che sia per
forza un game designer: l'anno precedente un team aveva come lead una
ragazza di programmazione. Il lead, precisa Andrea, è ``non tanto un capo
quanto un \textbf{rappresentante}'': la persona che scrive le mail di
aggiornamento, presenta i progetti e così via.

\subsection{I saluti}

Gli altri dubbi troveranno risposta nei prossimi incontri. Con gli studenti
di game design ci si vedrà ``lunedì non prossimo, ma quello dopo'' per
iniziare, e con Jack a novembre, quando cominceranno le sue lezioni; con
tutti gli altri ``ci si vedrà in giro''. ``Buon resto delle lezioni Growth.
Buon pranzo, anche.''
